\begin{figure}[h!]
    \centering
    \includegraphics[scale=0.20]{figures/graph_v1.png}
    \caption{Many point-lookups on the same key increase I/Os linearly. PLACEHOLDER. ADD GRAPH WITH DELETES IMPACT INSTEAD. MAYBE KEEP THAT}
    \label{fig:method}
\end{figure}

\Paragraph{LSM-based Key-Value Stores} 
To store all information Key-Value stores use the log-structured merge (LSM) tree.
LSM-trees store incoming data entries in main memory to facilitate writes, and periodically flush the entries as an immutable sorted run on durable storage.
Runs store entries in a specific type of file called 'SST files'.
During every compaction, the participating runs are sort-merged and invalidated entries are permanently deleted. Each compaction uses runs from consecutive levels.
LSM-trees are adopted by several modern systems including LevelDB~\cite{} and BigTable~\cite{} at Google, RocksDB~\cite{} at Facebook, X-Engine~\cite{} at Alibaba, Voldemort~\cite{} at LinkedIn, Dynamo~\cite{} at Amazon, Cassandra~\cite{}, HBase~\cite{}, and Accumulo~\cite{} at Apache, and bLSM~\cite{} and cLSM~\cite{} at Yahoo. 
MyRocks~\cite{} uses RocksDB as storage engine and SQLite4~\cite{} has experimented with LSM-trees in its storage layer, while columnar systems use LSM-style updates~\cite{}.

\Paragraph{Deletes in RocksDB}
It might not seem obvious at first but deletes are probably the most common operation in LSM trees and RocksDB. 
Of course deletes are used to invalidate specific entries, but they are also used during updates.
LSM trees use out-of-place updates which means that instead of updating the requested value, a new version is inserted in the higher levels. 
To complete this operation and avoid keeping older versions of values, LSM trees have to invalidate the older versions with a  tombstone. 
As a result, the delete operation accounts for almost half of every update operation. 
This, combined with the fact that deletes happen very frequently in many workloads, suggest that both range and point deletes majorly affect the performance such systems.
In this paper we will concentrate more on Range Deletes and ways to improve their performance.(MORE ON WHY RANGE DELETES IN SPECIFIC).
When a range delete is issued, a range tombstone is generated and added to the database. 
A point request initiates a process where all range tombstones are fetched from the disk and fragmented in memory. 
This enables the system to respond whether the requested key is deleted by such a range tombstone.


\Paragraph{The Problem} As previously mentioned, if a key is deleted by a range delete, we use one I/O to fetch the key and its sequence number and then ask the range delete aggregator to check if it is deleted. 
Thus, even if the key is deleted, we waste one I/O. 
If the key is requested for a second time, the same operation takes place, wasting more I/Os linearly.
This behavior can be seen in Figure 1.

\Paragraph{The Solution} 
Our method reduces disk accesses when a key deleted by a range delete is accessed multiple times.
It stores the deleted keys in a new key-value store which resides entirely in memory. 
This store contains deleted keys and gives no false-negative responses.
In this way, when a range deleted key is requested the first time, it is saved in this KV store. 
All future requests will first ask this KV store before accessing the disk to obtain the sequence number, thus eliminating I/Os.
In addition to that, we take advantage of the granularity that disk accesses happen.
Instead of retrieving and storing information in our KV store for only the requested key, we store information about all the keys in the same page.
In this way we can store more information and potentially avoid IOs on range deleted keys in the same page which may have not been requested yet.

